%==============================================================================
% Sjabloon poster bachproef
%==============================================================================
% Gebaseerd op document class `a0poster' door Gerlinde Kettl en Matthias Weiser
% Aangepast voor gebruik aan HOGENT door Jens Buysse en Bert Van Vreckem

\documentclass[a0,portrait]{hogent-poster}

% Info over de opleiding
\course{Bachelorproef}
\studyprogramme{Toegepaste Informatica}
\academicyear{2025-2026}
\institution{Hogeschool Gent, Valentin Vaerwyckweg 1, 9000 Gent}

% Info over de bachelorproef
\title{Detectie en analyse van Linux kernel-rootkits als educatieve leeromgeving}
\subtitle{Een veilige, representatieve en reproduceerbare labo-omgeving voor de cursus Cybersecurity Advanced}
\author{Elias Van Meerssche}
\email{elias.vanmeerssche@student.hogent.be}
\supervisor{Dhr. T. Clauwaert}
\cosupervisor{Dhr. B. Van Vreckem}

% Indien ingevuld, wordt deze informatie toegevoegd aan het einde van de
% abstract. Zet in commentaar als je dit niet wilt.
\specialisation{Systeem- en Netwerkbeheer}
\keywords{Linux, Kernel, Rootkit, eBPF, LKM, Cybersecurity}
\projectrepo{https://github.com/EliasVanMeerssche/Batcherproef-scripts}

\begin{document}
    
    \maketitle
    
    \begin{abstract}
        De kernel bevindt zich in het centrum van het besturingssysteem en beheert uiterst belangrijke taken[cite: 1]. Een rootkit is een soort malware die een aanvaller volledige controle geeft over het systeem en zich zeer diep in de kernel verbergt, wat detectie enorm bemoeilijkt[cite: 1]. Voor het vak Cybersecurity Advanced ontbrak het aan een veilige en reproduceerbare manier om deze kernel-corruptie te demonstreren. In dit onderzoek is een educatieve labo-omgeving ontworpen met Proof of Concept (PoC) virtuele machines (.ova-bestanden) om studenten Toegepaste Informatica te trainen in de detectie van zowel traditionele LKM- als moderne eBPF-rootkits[cite: 1].
    \end{abstract}
    
    \begin{multicols}{2} % This is how many columns your poster will be broken into
        
        \section{Probleemstelling \& Onderzoeksvragen}
        Omdat rootkits zo diep in het systeem zitten, is het moeilijk om dit op een eenvoudige en veilige manier te demonstreren tijdens de les[cite: 1]. Dit onderzoek beantwoordt de volgende hoofdvraag:
        
        \textbf{"Hoe kan een representatieve, veilige en reproduceerbare labo-omgeving worden ontworpen om studenten Toegepaste Informatica te trainen in de detectie en analyse van zowel traditionele (LKM) als moderne (eBPF) Linux-kernel-rootkits?"}[cite: 1]
        
        Dit wordt ondersteund door drie deelvragen rond[cite: 1]:
        \begin{itemize}
            \item \textbf{Theoretische basis:} Verschillen in hooking-technieken (ftrace vs. Kprobes).
            \item \textbf{Omgevingsconfiguratie:} Omzeilen van beveiligingen zoals SMEP, SMAP en KASLR.
            \item \textbf{Detectie \& Preventie:} Effectief gebruik van geheugenanalyse (Volatility) en runtime security (Falco, Tracee).
        \end{itemize}
        
        \section{Methodologie \& Opzet}
        Het onderzoek is opgedeeld in iteratieve fasen[cite: 1]:
        \begin{enumerate}
            \item \textbf{De klassieke dreiging (LKM):} Verdieping in een Loadable Kernel Module rootkit (\texttt{Singularity}) en het isoleren hiervan in een veilige \texttt{.ova} omgeving.
            \item \textbf{De moderne dreiging (eBPF):} Analyse van een Extended Berkeley Packet Filter rootkit (\texttt{bad-bpf}). In de labo-omgeving wordt de persistente inlaadmethode gegarandeerd via een \texttt{systemd}-service (\texttt{.system-health}).
            \item \textbf{Synthese en Preventie:} Het documenteren en testen van detectiestrategieën (zoals LiME en Volatility3) om dit te vormen tot een educatief stappenplan.
        \end{enumerate}
        Om geheugendumps succesvol te analyseren werd de kernelversie in de testomgeving vastgezet op \texttt{5.15.0-179}[cite: 1].
        
        \section{Achtergrond: LKM vs eBPF}
        \begin{center}
            \begin{tabular}{|l|p{14cm}|p{14cm}|}
                \hline
                \textbf{Kenmerk} & \textbf{LKM (Klassiek)} & \textbf{eBPF (Modern)} \\
                \hline
                \textbf{Inladen} & Fysiek modulebestand in besturingssysteem[cite: 1]. & Via \texttt{bpf()} systeemaanroep, in een virtuele machine[cite: 1]. \\
                \hline
                \textbf{Toegang} & Onbeperkt, risico op kernel panic bij fouten[cite: 1]. & Gelimiteerd, moet helperfuncties gebruiken (bv. \texttt{bpf\_probe\_write\_user})[cite: 1]. \\
                \hline
                \textbf{Hooking} & Direct overschrijven, vaak via \texttt{ftrace}[cite: 1]. & Via strikte inhaakpunten zoals Kprobes, Uprobes en Tracepoints[cite: 1]. \\
                \hline
            \end{tabular}
            \captionof{table}{Vergelijking van de offensieve technieken}
        \end{center}
        
        \section{Resultaten \& Detectie}
        Tijdens de testfase bleek dat de eBPF-rootkit (\texttt{ebpfkit}) het inladen van LKM's blokkeert om geheugenforensisch onderzoek via tools zoals LiME tegen te gaan[cite: 1]. Om deze reden wordt een hypervisor-dump (\texttt{VBoxManage dumpvmcore}) gebruikt als betrouwbaar alternatief[cite: 1].
        
        Bij analyse met \textbf{Volatility 3} (\texttt{linux.ebpf.EBPF}) worden de kwaadaardige BPF-programma's in het geheugen zichtbaar[cite: 1]:
        
        \begin{verbatim}
            Address        Name             Tag              Type
            0xccbb80051000 kretprobe_64_s   72a675d136f86ad3 BPF_PROG_TYPE_KPROBE
            0xccbb806fb000 xdp_ingress_del  b489bab0c6d5985c BPF_PROG_TYPE_XDP
            0xccbb8322b000 egress_dispatch  f7ef2bda3a4006f0 BPF_PROG_TYPE_SCHED_CLS
        \end{verbatim}
        
        Hoewel het type \texttt{BPF\_PROG\_TYPE\_CGROUP\_SKB} legitiem is, toont de aanwezigheid van \texttt{KPROBE}, \texttt{XDP} en \texttt{SCHED\_CLS} de verborgen rootkit direct aan[cite: 1]. Ook commando's zoals \texttt{dmesg} en tools zoals \texttt{Falco} genereren waarschuwingen wanneer de \texttt{bpf\_probe\_write\_user} helperfunctie wordt gebruikt[cite: 1].
        
        \section{Conclusies}
        De primaire doelstelling is succesvol bereikt: er is een werkende Proof of Concept (PoC) opgeleverd in de vorm van \texttt{.ova}-bestanden (met de LKM rootkit \texttt{Singularity} en eBPF rootkit \texttt{bad-bpf})[cite: 1]. Hierdoor is het voorheen abstracte en lastig te doceren gevaar van rootkits vertaald naar een concrete, begrijpelijke labo-oefening. Studenten kunnen nu op een veilige manier de theorie achter kernel-corruptie oefenen en leren werken met detectietools zoals Volatility, Falco en Tracee[cite: 1].
        
        \section{Beperkingen \& Toekomstig onderzoek}
        \textbf{Beperkingen van dit onderzoek:} Het daadwerkelijk uitbuiten van echte kwetsbaarheden (privilege escalation) is in de methodologie weggelaten om de complexiteit te verlagen en de focus op detectie te behouden[cite: 1]. Daarnaast wordt eBPF-persistentie gesimuleerd met een \texttt{systemd}-service in plaats van via geavanceerde bootkit-mechanismen[cite: 1].
        
        \textbf{Suggesties voor verder onderzoek:} Het is waardevol om de geautomatiseerde besmetting via echte kwetsbaarheden verder uit te werken[cite: 1]. Daarnaast kan de impact van command \& control netwerkmanipulatie via XDP (Express Data Path) hooks grondiger geanalyseerd worden binnen de testomgeving[cite: 1].
        
    \end{multicols}
\end{document}